\par\medskip\noindent
\underline{Corolario 2.10.} (a) Sea
$\varphi:\underline A\longrightarrow\underline B$ un funtor liso de
grupoides cofibrados sobre $\underline C_\Lambda$. Sea $\xi$ un objeto
formal versal de $\underline A$, y sea $\eta$ un objeto formal
mínimamente versal de $\underline B$. Si $\underline B$ es homogéneo,
entonces, para todo morfismo
$\beta:\eta\longrightarrow\varphi(\xi)$ en
$\widehat{\underline B}$, el anillo $P(\xi)$ es un anillo de series
formales sobre $P(\eta)$ mediante
\[
 P(\beta):P(\eta)\longrightarrow P(\varphi(\xi))=P(\xi).
\]

(b) Sea $\underline A$ un grupoide cofibrado homogéneo sobre
$\underline C_\Lambda$, y sean $\xi$ y $\eta$ objetos formales versales
de $\underline A$. Si $\eta$ es minimal, entonces, para todo morfismo
$\alpha:\eta\longrightarrow\xi$ en $\widehat{\underline A}$, el
anillo $P(\xi)$ es un anillo de series formales sobre $P(\eta)$ mediante
\[
 P(\alpha):P(\eta)\longrightarrow P(\xi).
\]

\underline{Demostración.} (a) Un morfismo
$\beta:\eta\longrightarrow\varphi(\xi)$ induce un diagrama conmutativo
\[
\begin{tikzcd}[row sep=huge,column sep=huge]
 \underline h_{P(\xi)} \arrow[r,"{\widetilde\beta}"]
                           \arrow[d,"{\widetilde\xi}"']
 & \underline h_{P(\eta)} \arrow[d,"{\widetilde\eta}"]\\
 \underline A \arrow[r,"{\varphi}"'] & \underline B.
\end{tikzcd}
\]
Como $\widetilde\xi$ y $\varphi$ son lisos, el compuesto
$\widetilde\eta\circ\widetilde\beta
=\varphi\circ\widetilde\xi$ es liso. Por otra parte,

%%%%%%%%%%  scan idx 71  |  folio 60  %%%%%%%%%%

\[
 \widetilde\eta(k[\varepsilon]):
 \underline h_{P(\eta)}(k[\varepsilon])
 \longrightarrow[\underline B(k[\varepsilon])]
\]
es un isomorfismo. La Proposición 2.9(a) muestra entonces que
$\widetilde\beta$ es liso, lo que significa precisamente que $P(\xi)$
es un anillo de series formales sobre $P(\eta)$.

(b) Apliquemos (a) con $\underline A=\underline B$.
\hfill$\square$
